\documentclass[multi=tikzpicture,border=9pt]{standalone}
\usepackage[UTF8]{ctex}
\usepackage{amsmath,amssymb}
\usepackage{tikz}
\usetikzlibrary{arrows.meta,calc,patterns,decorations.pathmorphing}
\definecolor{main}{HTML}{17365D}
\definecolor{accent}{HTML}{AC3931}
\tikzset{every picture/.style={font=\small,line width=.75pt},
>={Stealth[length=2mm]},arr/.style={->},
box/.style={draw=main,rounded corners=2pt,fill=main!4,text width=3cm,minimum height=1cm,align=center},
note/.style={align=center,text width=9cm}}
\begin{document}
\begin{tikzpicture}
\draw[arr] (-3,0)--(4,0) node[right]{分子轴};\foreach \x/\lab in {-1/-q,0/2q,1/-q}{\fill[accent] (\x,0) circle (3pt);\node[below] at (\x,-.1){$\lab$};}
\draw[<->] (-1,-.8)--(0,-.8) node[midway,below]{$l$};\draw[<->] (0,-.8)--(1,-.8) node[midway,below]{$l$};\draw[arr,main] (0,0)--(40:3.8) node[midway,above]{$r$};\fill[main] (40:3.8) circle (2pt);\draw[arr] (.7,0) arc (0:40:.7) node[midway,right]{$\theta$};\node[note] at (.5,-2){$Q_{\rm net}=0$，$\mathbf p=0$，四极项不为零};
\end{tikzpicture}
\begin{tikzpicture}
\node[box] (n0) at (0,0) {单极项\\$\sum q_i=0$};
\node[box] (n1) at (5,0) {偶极项\\$\sum q_i\mathbf r_i=0$};
\node[box] (n2) at (5,-2.3) {四极电势\\$V\propto r^{-3}$};
\node[box] (n3) at (0,-2.3) {取完整梯度\\$\mathbf E=-\boldsymbol\nabla V$};
\draw[arr] (n0)--(n1);\draw[arr] (n1)--(n2);\draw[arr] (n2)--(n3);
\node[note] at (2.5,-4.1) {左右对称使奇数阶项消失；余项从 $l^4/r^5$ 起};
\end{tikzpicture}
\begin{tikzpicture}
\draw[dashed] (-3,0)--(3,0);\draw[dashed] (0,-3)--(0,3);\node[right] at (3,0){分子轴};
\draw[arr,main] (1.7000,0.0000)--(1.2800,0.0000);
\draw[arr,main] (1.5706,0.6506)--(1.2835,0.3440);
\draw[arr,main] (1.2021,1.2021)--(1.3349,0.8036);
\draw[arr,main] (0.6506,1.5706)--(1.0545,1.6857);
\draw[arr,main] (0.0000,1.7000)--(0.0000,2.1200);
\draw[arr,main] (-0.6506,1.5706)--(-1.0545,1.6857);
\draw[arr,main] (-1.2021,1.2021)--(-1.3349,0.8036);
\draw[arr,main] (-1.5706,0.6506)--(-1.2835,0.3440);
\draw[arr,main] (-1.7000,0.0000)--(-1.2800,0.0000);
\draw[arr,main] (-1.5706,-0.6506)--(-1.2835,-0.3440);
\draw[arr,main] (-1.2021,-1.2021)--(-1.3349,-0.8036);
\draw[arr,main] (-0.6506,-1.5706)--(-1.0545,-1.6857);
\draw[arr,main] (-0.0000,-1.7000)--(-0.0000,-2.1200);
\draw[arr,main] (0.6506,-1.5706)--(1.0545,-1.6857);
\draw[arr,main] (1.2021,-1.2021)--(1.3349,-0.8036);
\draw[arr,main] (1.5706,-0.6506)--(1.2835,-0.3440);
\draw[arr,main] (2.6000,0.0000)--(2.1800,0.0000);
\draw[arr,main] (2.4021,0.9950)--(2.1150,0.6884);
\draw[arr,main] (1.8385,1.8385)--(1.9713,1.4400);
\draw[arr,main] (0.9950,2.4021)--(1.3989,2.5172);
\draw[arr,main] (0.0000,2.6000)--(0.0000,3.0200);
\draw[arr,main] (-0.9950,2.4021)--(-1.3989,2.5172);
\draw[arr,main] (-1.8385,1.8385)--(-1.9713,1.4400);
\draw[arr,main] (-2.4021,0.9950)--(-2.1150,0.6884);
\draw[arr,main] (-2.6000,0.0000)--(-2.1800,0.0000);
\draw[arr,main] (-2.4021,-0.9950)--(-2.1150,-0.6884);
\draw[arr,main] (-1.8385,-1.8385)--(-1.9713,-1.4400);
\draw[arr,main] (-0.9950,-2.4021)--(-1.3989,-2.5172);
\draw[arr,main] (-0.0000,-2.6000)--(-0.0000,-3.0200);
\draw[arr,main] (0.9950,-2.4021)--(1.3989,-2.5172);
\draw[arr,main] (1.8385,-1.8385)--(1.9713,-1.4400);
\draw[arr,main] (2.4021,-0.9950)--(2.1150,-0.6884);
\fill[accent] (-.45,0) circle (2pt);\fill[accent] (0,0) circle (2pt);\fill[accent] (.45,0) circle (2pt);\node[note] at (0,-4){箭头仅示方向，长度统一；轴向向内、中垂面向外};
\end{tikzpicture}
\begin{tikzpicture}
\draw[dashed] (-2,0)--(7,0);\foreach \off in {0,5}{\foreach \x/\lab in {-1/-q,0/2q,1/-q}{\fill[main] (\off+\x,0) circle (3pt);\node[above] at (\off+\x,.15){$\lab$};}}
\draw[<->] (0,-1)--(5,-1) node[midway,below]{$R\gg l$};\draw[arr,accent,very thick] (-1.3,0)--(-2.3,0);\draw[arr,accent,very thick] (6.3,0)--(7.3,0);
\node[note] at (2.5,-2.6){$U_{12}=2qV(R)-qV(R-l)-qV(R+l)$\\$U_{12}\simeq24kq^2l^4/R^5>0$，共轴排列主阶排斥};
\end{tikzpicture}
\begin{tikzpicture}
\node[box] (n0) at (0,0) {四极电势\\$V\propto r^{-3}$};
\node[box] (n1) at (5,0) {单分子电场\\$E\propto r^{-4}$};
\node[box] (n2) at (5,-2.3) {两分子作用能\\$U_{12}\propto R^{-5}$};
\node[box] (n3) at (0,-2.3) {两分子合力\\$F\propto R^{-6}$};
\draw[arr] (n0)--(n1);\draw[arr] (n1)--(n2);\draw[arr] (n2)--(n3);
\node[note] at (2.5,-4.1) {第二分子的中性与对称性再次抵消低阶项；固定共轴方向};
\end{tikzpicture}
\end{document}
